% Part: first-order-logic
% Chapter: completeness
% Section: henkin-expansion

\documentclass[../../../include/open-logic-section]{subfiles}

\begin{document}

\olfileid{fol}{com}{hen}

\olsection{توسيع هنكن}

\begin{explain}
من مطالب برهان الاكتمال أن النموذج المبني من مجموعة متسقة
كاملة~${\OLId{greek-variable}{Gamma}}$ يرضي ال!!{formula}s المكممة كلها في~${\OLId{greek-variable}{Gamma}}$.
ولتحصيل هذا الضمان نأخذ بطريقة ليون هنكن: نزيد اللغة
!!{constant}s لا نهاية لها، ونلحق، لكل !!{formula} لها
!!{variable} حر واحد $!A({\OLId{latin-ordinary}{x}})$، صيغة على النحو
\iftag{prvEx}
      {$\lexists[{\OLId{latin-ordinary}{x}}][!A({\OLId{latin-ordinary}{x}})] \lif !A({\OLId{latin-ordinary}{c}})$}
      {$\lnot\lforall[{\OLId{latin-ordinary}{x}}][!A({\OLId{latin-ordinary}{x}})] \lif \lnot !A({\OLId{latin-ordinary}{c}})$},
على أن ${\OLId{latin-ordinary}{c}}$ من ال!!{constant}s المستحدثة. فإذا بنينا ال!!{structure}
المرضية لـ~${\OLId{greek-variable}{Gamma}}$، كفلت هذه الزيادة لكل
\iftag{prvEx}
{جملة وجودية صادقة شاهدًا}
{جملة كلية كاذبة مثالًا ناقضًا}
من بين الرموز الثابتة الجديدة.
\end{explain}

\begin{prop}
\ollabel{prop:lang-exp}
إذا كانت ${\OLId{greek-variable}{Gamma}}$ متسقة في $\Lang L$، وكانت $\Lang L'$ ناتجة من
$\Lang L$ بإضافة مجموعة !!a{denumerable} من ال!!{constant}s الجديدة
$\Obj d_{\OLMathNumeral{0}}$ و$\Obj d_{\OLMathNumeral{1}}$ و\dots، فإن ${\OLId{greek-variable}{Gamma}}$ متسقة في~$\Lang L'$.
\end{prop}

\begin{defn}[المجموعة المشبعة]
المجموعة ${\OLId{greek-variable}{Gamma}}$ من !!{formula}s اللغة $\Lang {L}$ تسمى
\emph{مشبعة} إذا وفقط إذا كان لكل !!{formula}~$!A({\OLId{latin-ordinary}{x}}) \in \Frm[{\OLId{latin-looped}{L}}]$
ليس لها إلا !!{variable} حر واحد~${\OLId{latin-ordinary}{x}}$، !!a{constant}~${\OLId{latin-ordinary}{c}} \in \Lang{L}$ يحقق
\iftag{prvEx}
      {$\lexists[{\OLId{latin-ordinary}{x}}][!A({\OLId{latin-ordinary}{x}})] \lif !A({\OLId{latin-ordinary}{c}}) \in {\OLId{greek-variable}{Gamma}}$}
      {$\lnot\lforall[{\OLId{latin-ordinary}{x}}][!A({\OLId{latin-ordinary}{x}})] \lif \lnot !A({\OLId{latin-ordinary}{c}}) \in {\OLId{greek-variable}{Gamma}}$}.
\end{defn}

ولأجل البرهان الآتي نضع التعريف التالي.

\begin{defn}
\ollabel{defn:henkin-exp}
لتكن $\Lang L'$ اللغة الموصوفة في \olref{prop:lang-exp}. نختار تعدادًا
$!A_{\OLMathNumeral{0}}({\OLId{latin-ordinary}{x}}_{\OLMathNumeral{0}})$ و$!A_{\OLMathNumeral{1}}({\OLId{latin-ordinary}{x}}_{\OLMathNumeral{1}})$ و\dots يستوعب ال!!{formula}s~$!A_{\OLId{latin-ordinary}{i}}({\OLId{latin-ordinary}{x}}_{\OLId{latin-ordinary}{i}})$
في~$\Lang L'$ التي لها متغير حر واحد (${\OLId{latin-ordinary}{x}}_{\OLId{latin-ordinary}{i}}$). ثم نبني
ال!!{sentence}s~$!D_{\OLId{latin-ordinary}{n}}$ استقراءً على~${\OLId{latin-ordinary}{n}}$.

فنختار ${\OLId{latin-ordinary}{c}}_{\OLMathNumeral{0}}$ أول !!{constant} في قائمة $\Obj d_{\OLId{latin-ordinary}{i}}$ المضافة
إلى~$\Lang{L}$ مما لا يقع في~$!A_{\OLMathNumeral{0}}({\OLId{latin-ordinary}{x}}_{\OLMathNumeral{0}})$. وإذا سبق تعريف $!D_{\OLMathNumeral{0}}$
و\dots و~$!D_{{\OLId{latin-ordinary}{n}}-{\OLMathNumeral{1}}}$، اخترنا ${\OLId{latin-ordinary}{c}}_{\OLId{latin-ordinary}{n}}$ أول ال!!{constant}s الجديدة
~$\Obj d_{\OLId{latin-ordinary}{i}}$ مما لا يقع في $!D_{\OLMathNumeral{0}}$ و\dots و~$!D_{{\OLId{latin-ordinary}{n}}-{\OLMathNumeral{1}}}$، ولا
في~$!A_{\OLId{latin-ordinary}{n}}({\OLId{latin-ordinary}{x}}_{\OLId{latin-ordinary}{n}})$.

ونعرّف $!D_{{\OLId{latin-ordinary}{n}}}$ بأنها ال!!{sentence} \iftag{prvEx}
{$\lexists[{\OLId{latin-ordinary}{x}}_{{\OLId{latin-ordinary}{n}}}][!A_{{\OLId{latin-ordinary}{n}}}({\OLId{latin-ordinary}{x}}_{{\OLId{latin-ordinary}{n}}})] \lif
  !A_{{\OLId{latin-ordinary}{n}}}({\OLId{latin-ordinary}{c}}_{{\OLId{latin-ordinary}{n}}})$}{$\lnot\lforall[{\OLId{latin-ordinary}{x}}_{{\OLId{latin-ordinary}{n}}}][!A_{{\OLId{latin-ordinary}{n}}}({\OLId{latin-ordinary}{x}}_{{\OLId{latin-ordinary}{n}}})] \lif \lnot
  !A_{{\OLId{latin-ordinary}{n}}}({\OLId{latin-ordinary}{c}}_{{\OLId{latin-ordinary}{n}}})$}.
\end{defn}

\begin{lem}
\ollabel{lem:henkin}
يمكن توسيع كل مجموعة متسقة~${\OLId{greek-variable}{Gamma}}$ إلى مجموعة متسقة مشبعة~${\OLId{greek-variable}{Gamma}}'$.
\end{lem}

\begin{proof}
ابدأ بمجموعة جمل متسقة~${\OLId{greek-variable}{Gamma}}$ في لغة~$\Lang{L}$، وزد اللغة مجموعة
!!a{denumerable} من ال!!{constant}s الجديدة، ولتكن اللغة الناتجة
~$\Lang{L'}$. تضمن \olref{prop:lang-exp} بقاء اتساق ${\OLId{greek-variable}{Gamma}}$ بعد
الزيادة. وخذ $!D_{\OLId{latin-ordinary}{i}}$ من \olref{defn:henkin-exp}، وعرّف
\begin{align*}
{\OLId{greek-variable}{Gamma}}_{\OLMathNumeral{0}} & = {\OLId{greek-variable}{Gamma}} \\
{\OLId{greek-variable}{Gamma}}_{{\OLId{latin-ordinary}{n}}+{\OLMathNumeral{1}}} & = {\OLId{greek-variable}{Gamma}}_{\OLId{latin-ordinary}{n}} \cup \{!D_{\OLId{latin-ordinary}{n}} \}
\end{align*}
فيكون ${\OLId{greek-variable}{Gamma}}_{{\OLId{latin-ordinary}{n}}+{\OLMathNumeral{1}}} = {\OLId{greek-variable}{Gamma}} \cup \{ !D_{\OLMathNumeral{0}}, \dots, !D_{\OLId{latin-ordinary}{n}} \}$. ثم ضع
${\OLId{greek-variable}{Gamma}}' = \bigcup_{{\OLId{latin-ordinary}{n}}} {\OLId{greek-variable}{Gamma}}_{\OLId{latin-ordinary}{n}}$؛ وإشباع ${\OLId{greek-variable}{Gamma}}'$ بيّن من بنائها.

إن عدم اتساق ${\OLId{greek-variable}{Gamma}}'$ يستلزم عدم اتساق ${\OLId{greek-variable}{Gamma}}_{\OLId{latin-ordinary}{n}}$ عند بعض ${\OLId{latin-ordinary}{n}}$
(بيّن وجه ذلك تمرينًا). فإثبات اتساق ${\OLId{greek-variable}{Gamma}}'$ مردود إلى الاستقراء
على~${\OLId{latin-ordinary}{n}}$ لإثبات اتساق كل~${\OLId{greek-variable}{Gamma}}_{\OLId{latin-ordinary}{n}}$.

أما الأساس، فاتساق ${\OLId{greek-variable}{Gamma}}_{\OLMathNumeral{0}} = {\OLId{greek-variable}{Gamma}}$ هو الفرض نفسه. وأما الانتقال،
فلنفرض اتساق ${\OLId{greek-variable}{Gamma}}_{{\OLId{latin-ordinary}{n}}}$ وعدم اتساق
${\OLId{greek-variable}{Gamma}}_{{\OLId{latin-ordinary}{n}}+{\OLMathNumeral{1}}} = {\OLId{greek-variable}{Gamma}}_{\OLId{latin-ordinary}{n}} \cup \{!D_{\OLId{latin-ordinary}{n}}\}$. وقد عرّفنا $!D_{\OLId{latin-ordinary}{n}}$~بأنها
\iftag{prvEx}
{$\lexists[{\OLId{latin-ordinary}{x}}_{{\OLId{latin-ordinary}{n}}}][!A_{{\OLId{latin-ordinary}{n}}}({\OLId{latin-ordinary}{x}}_{\OLId{latin-ordinary}{n}})] \lif !A_{{\OLId{latin-ordinary}{n}}}({\OLId{latin-ordinary}{c}}_{{\OLId{latin-ordinary}{n}}})$}
{$\lnot\lforall[{\OLId{latin-ordinary}{x}}_{{\OLId{latin-ordinary}{n}}}][!A_{{\OLId{latin-ordinary}{n}}}({\OLId{latin-ordinary}{x}}_{\OLId{latin-ordinary}{n}})] \lif \lnot !A_{{\OLId{latin-ordinary}{n}}}({\OLId{latin-ordinary}{c}}_{{\OLId{latin-ordinary}{n}}})$},
وصيغة $!A_{\OLId{latin-ordinary}{n}}({\OLId{latin-ordinary}{x}}_{\OLId{latin-ordinary}{n}})$ هي !!a{formula} من $\Lang{L'}$ متغيرها الحر
الوحيد~${\OLId{latin-ordinary}{x}}_{\OLId{latin-ordinary}{n}}$. وطريقة اختيار~${\OLId{latin-ordinary}{c}}_{\OLId{latin-ordinary}{n}}$ في
\olref{defn:henkin-exp} تكفل أن ${\OLId{latin-ordinary}{c}}_{\OLId{latin-ordinary}{n}}$ لا يقع في~$!A_{\OLId{latin-ordinary}{n}}({\OLId{latin-ordinary}{x}}_{\OLId{latin-ordinary}{n}})$ ولا في~${\OLId{greek-variable}{Gamma}}_{\OLId{latin-ordinary}{n}}$.

فمن عدم اتساق ${\OLId{greek-variable}{Gamma}}_{\OLId{latin-ordinary}{n}} \cup \{!D_{\OLId{latin-ordinary}{n}}\}$ يلزم ${\OLId{greek-variable}{Gamma}}_{\OLId{latin-ordinary}{n}}
\Proves \lnot !D_{\OLId{latin-ordinary}{n}}$، ويلزم من ذلك هذان الحكمان:
\iftag{prvEx}{
\[
{\OLId{greek-variable}{Gamma}}_{\OLId{latin-ordinary}{n}} \Proves \lexists[{\OLId{latin-ordinary}{x}}_{\OLId{latin-ordinary}{n}}][!A_{\OLId{latin-ordinary}{n}}({\OLId{latin-ordinary}{x}}_{\OLId{latin-ordinary}{n}})]
\qquad
{\OLId{greek-variable}{Gamma}}_{\OLId{latin-ordinary}{n}} \Proves \lnot !A_{\OLId{latin-ordinary}{n}}({\OLId{latin-ordinary}{c}}_{\OLId{latin-ordinary}{n}})
\]}{
\[
{\OLId{greek-variable}{Gamma}}_{\OLId{latin-ordinary}{n}} \Proves \lnot\lforall[{\OLId{latin-ordinary}{x}}_{\OLId{latin-ordinary}{n}}][!A_{\OLId{latin-ordinary}{n}}({\OLId{latin-ordinary}{x}}_{\OLId{latin-ordinary}{n}})]
\qquad
{\OLId{greek-variable}{Gamma}}_{\OLId{latin-ordinary}{n}} \Proves !A_{\OLId{latin-ordinary}{n}}({\OLId{latin-ordinary}{c}}_{\OLId{latin-ordinary}{n}})
\]}
ولأن ${\OLId{latin-ordinary}{c}}_{\OLId{latin-ordinary}{n}}$ لا يقع في ${\OLId{greek-variable}{Gamma}}_{\OLId{latin-ordinary}{n}}$ ولا في~$!A_{\OLId{latin-ordinary}{n}}({\OLId{latin-ordinary}{x}}_{\OLId{latin-ordinary}{n}})$، يصح تطبيق
\tagrefs{prfAX/{fol:axd:qpr:thm:strong-generalization},prfSC/{fol:seq:qpr:thm:strong-generalization},prfND/{fol:ntd:qpr:thm:strong-generalization},prfTab/{fol:tab:qpr:thm:strong-generalization}}.
ومن \iftag{prvEx}
{${\OLId{greek-variable}{Gamma}}_{\OLId{latin-ordinary}{n}} \Proves \lnot !A_{\OLId{latin-ordinary}{n}}({\OLId{latin-ordinary}{c}}_{\OLId{latin-ordinary}{n}})$}
{${\OLId{greek-variable}{Gamma}}_{\OLId{latin-ordinary}{n}} \Proves !A_{\OLId{latin-ordinary}{n}}({\OLId{latin-ordinary}{c}}_{\OLId{latin-ordinary}{n}})$},
نحصل على
\iftag{prvEx}
{${\OLId{greek-variable}{Gamma}}_{\OLId{latin-ordinary}{n}} \Proves \lforall[{\OLId{latin-ordinary}{x}}_{\OLId{latin-ordinary}{n}}][\lnot !A_{\OLId{latin-ordinary}{n}}({\OLId{latin-ordinary}{x}}_{\OLId{latin-ordinary}{n}})]$}
{${\OLId{greek-variable}{Gamma}}_{\OLId{latin-ordinary}{n}} \Proves \lforall[{\OLId{latin-ordinary}{x}}_{\OLId{latin-ordinary}{n}}][!A_{\OLId{latin-ordinary}{n}}({\OLId{latin-ordinary}{x}}_{\OLId{latin-ordinary}{n}})]$}.
فيجتمع الحكمان
\iftag{prvEx}
{${\OLId{greek-variable}{Gamma}}_{\OLId{latin-ordinary}{n}} \Proves \lexists[{\OLId{latin-ordinary}{x}}_{\OLId{latin-ordinary}{n}}][!A_{\OLId{latin-ordinary}{n}}({\OLId{latin-ordinary}{x}}_{\OLId{latin-ordinary}{n}})]$ و
${\OLId{greek-variable}{Gamma}}_{\OLId{latin-ordinary}{n}} \Proves \lforall[{\OLId{latin-ordinary}{x}}_{\OLId{latin-ordinary}{n}}][\lnot !A_{\OLId{latin-ordinary}{n}}({\OLId{latin-ordinary}{x}}_{\OLId{latin-ordinary}{n}})]$}
{${\OLId{greek-variable}{Gamma}}_{\OLId{latin-ordinary}{n}} \Proves \lnot\lforall[{\OLId{latin-ordinary}{x}}_{\OLId{latin-ordinary}{n}}][!A_{\OLId{latin-ordinary}{n}}({\OLId{latin-ordinary}{x}}_{\OLId{latin-ordinary}{n}})]$ و
${\OLId{greek-variable}{Gamma}}_{\OLId{latin-ordinary}{n}} \Proves \lforall[{\OLId{latin-ordinary}{x}}_{\OLId{latin-ordinary}{n}}][!A_{\OLId{latin-ordinary}{n}}({\OLId{latin-ordinary}{x}}_{\OLId{latin-ordinary}{n}})]$},
ومن اجتماعهما يلزم عدم اتساق ${\OLId{greek-variable}{Gamma}}_{\OLId{latin-ordinary}{n}}$ نفسها.
\iftag{prvEx}
{(لاحظ أن
\iftag{prvAll}
{$\lforall[{\OLId{latin-ordinary}{x}}_{\OLId{latin-ordinary}{n}}][\lnot !A_{\OLId{latin-ordinary}{n}}({\OLId{latin-ordinary}{x}}_{\OLId{latin-ordinary}{n}})] \Proves
  \lnot\lexists[{\OLId{latin-ordinary}{x}}_{\OLId{latin-ordinary}{n}}][!A_{\OLId{latin-ordinary}{n}}({\OLId{latin-ordinary}{x}}_{\OLId{latin-ordinary}{n}})]$}
{$\lforall[{\OLId{latin-ordinary}{x}}_{\OLId{latin-ordinary}{n}}][\lnot !A_{\OLId{latin-ordinary}{n}}({\OLId{latin-ordinary}{x}}_{\OLId{latin-ordinary}{n}})]$ تُعرّف بأنها
  $\lnot\lexists[{\OLId{latin-ordinary}{x}}_{\OLId{latin-ordinary}{n}}][\lnot\lnot !A_{\OLId{latin-ordinary}{n}}({\OLId{latin-ordinary}{x}}_{\OLId{latin-ordinary}{n}})]$ و
  $\lnot\lexists[{\OLId{latin-ordinary}{x}}_{\OLId{latin-ordinary}{n}}][\lnot\lnot !A_{\OLId{latin-ordinary}{n}}({\OLId{latin-ordinary}{x}}_{\OLId{latin-ordinary}{n}})] \Proves
  \lnot\lexists[{\OLId{latin-ordinary}{x}}_{\OLId{latin-ordinary}{n}}][!A_{\OLId{latin-ordinary}{n}}({\OLId{latin-ordinary}{x}}_{\OLId{latin-ordinary}{n}})]$}.)}{}
وهذا خلاف اتساق ${\OLId{greek-variable}{Gamma}}_{\OLId{latin-ordinary}{n}}$ المفروض. فثبت اتساق
${\OLId{greek-variable}{Gamma}}_{\OLId{latin-ordinary}{n}} \cup \{ !D_{\OLId{latin-ordinary}{n}}\}$، وتم الاستقراء.
\end{proof}

\begin{explain}
ننتقل إلى المجموعات التي يجتمع فيها الاتساق و\emph{الاكتمال} والإشباع.
سنثبت أن المجموعة من هذا القبيل
\iftag{prvAll}{تضم !!{sentence} مكممة كليًا إذا وفقط إذا ضمت
  جميع حالاتها}{}\iftag{defAll,defEx}{}{، وأنها }\iftag{prvEx}{تضم
  !!{sentence} مكممة وجوديًا إذا وفقط إذا ضمت حالة واحدة منها
  على الأقل}{}. بهذه الخاصية نثبت أن ال!!{structure} المنشأة من
مجموعة كاملة متسقة مشبعة ترضي جميع جمل المجموعة المكممة.
\end{explain}

\begin{prop}\ollabel{prop:saturated-instances}
افترض أن ${\OLId{greek-variable}{Gamma}}$ كاملة ومتسقة ومشبعة.
\begin{tagenumerate}{prvEx,prvAll}
\tagitem{prvEx}{$\lexists[{\OLId{latin-ordinary}{x}}][!A({\OLId{latin-ordinary}{x}})] \in {\OLId{greek-variable}{Gamma}}$ إذا وفقط إذا كان
  $!A({\OLId{latin-ordinary}{t}}) \in {\OLId{greek-variable}{Gamma}}$ لحد مغلق واحد~${\OLId{latin-ordinary}{t}}$ على الأقل.}{}
\tagitem{prvAll}{$\lforall[{\OLId{latin-ordinary}{x}}][!A({\OLId{latin-ordinary}{x}})] \in {\OLId{greek-variable}{Gamma}}$ إذا وفقط إذا كان
  $!A({\OLId{latin-ordinary}{t}}) \in {\OLId{greek-variable}{Gamma}}$ لكل حد مغلق~${\OLId{latin-ordinary}{t}}$.}{}
\end{tagenumerate}
\end{prop}

\begin{proof}
  \begin{tagenumerate}{prvEx,prvAll}
  \tagitem{prvEx}{%
    \iftag{probEx}{تمرين.}{أما إذا كان $\lexists[{\OLId{latin-ordinary}{x}}][!A({\OLId{latin-ordinary}{x}})]
      \in {\OLId{greek-variable}{Gamma}}$، فإن إشباع ${\OLId{greek-variable}{Gamma}}$ يضمن
      $(\lexists[{\OLId{latin-ordinary}{x}}][!A({\OLId{latin-ordinary}{x}})] \lif !A({\OLId{latin-ordinary}{c}})) \in {\OLId{greek-variable}{Gamma}}$ لبعض
      ال!!{constant}~${\OLId{latin-ordinary}{c}}$. وبحسب
      \tagrefs{prfAX/{fol:axd:ppr:prop:provability-lif},prfSC/{fol:seq:ppr:prop:provability-lif},prfND/{fol:ntd:ppr:prop:provability-lif},prfTab/{fol:tab:ppr:prop:provability-lif}}، البند~(\OLClassicalDigits{1})،
      و\olref[ccs]{prop:ccs}\olref[ccs]{prop:ccs-prov-in}، نحصل على
      $!A({\OLId{latin-ordinary}{c}}) \in {\OLId{greek-variable}{Gamma}}$.

    وأما العكس فلا يتوقف على الإشباع. فمتى كان $!A({\OLId{latin-ordinary}{t}}) \in {\OLId{greek-variable}{Gamma}}$،
    ثبت ${\OLId{greek-variable}{Gamma}} \Proves \lexists[{\OLId{latin-ordinary}{x}}][!A({\OLId{latin-ordinary}{x}})]$ بمقتضى
      \tagrefs{prfAX/{fol:axd:qpr:prop:provability-quantifiers},prfSC/{fol:seq:qpr:prop:provability-quantifiers},prfND/{fol:ntd:qpr:prop:provability-quantifiers},prfTab/{fol:tab:qpr:prop:provability-quantifiers}}، البند~(\OLClassicalDigits{1}). وبحسب
    \olref[ccs]{prop:ccs}\olref[ccs]{prop:ccs-prov-in} نحصل على
    $\lexists[{\OLId{latin-ordinary}{x}}][!A({\OLId{latin-ordinary}{x}})] \in {\OLId{greek-variable}{Gamma}}$.}}{}
  \tagitem{prvAll}{%
    \iftag{probAll}{تمرين.}{ليكن $!A({\OLId{latin-ordinary}{t}}) \in {\OLId{greek-variable}{Gamma}}$ لكل حد
      مغلق~${\OLId{latin-ordinary}{t}}$. فلو كان
      $\lforall[{\OLId{latin-ordinary}{x}}][!A({\OLId{latin-ordinary}{x}})] \notin {\OLId{greek-variable}{Gamma}}$، لأوجب اكتمال ${\OLId{greek-variable}{Gamma}}$
      أن $\lnot\lforall[{\OLId{latin-ordinary}{x}}][!A({\OLId{latin-ordinary}{x}})] \in {\OLId{greek-variable}{Gamma}}$. وبالإشباع نحصل على
      \iftag{prvEx}{$(\lexists[{\OLId{latin-ordinary}{x}}][\lnot !A({\OLId{latin-ordinary}{x}})] \lif \lnot !A({\OLId{latin-ordinary}{c}})) \in
        {\OLId{greek-variable}{Gamma}}$}{$(\lnot\lforall[{\OLId{latin-ordinary}{x}}][!A({\OLId{latin-ordinary}{x}})] \lif \lnot !A({\OLId{latin-ordinary}{c}})) \in
        {\OLId{greek-variable}{Gamma}}$} لبعض ال!!{constant}~${\OLId{latin-ordinary}{c}}$. ثم إن ${\OLId{latin-ordinary}{c}}$
      حد مغلق، فيقتضي الفرض $!A({\OLId{latin-ordinary}{c}}) \in {\OLId{greek-variable}{Gamma}}$؛ وحينئذ لا تكون ${\OLId{greek-variable}{Gamma}}$
      متسقة. (تمرين: أقم ال!!{derivation} المبين أن
      \[
      \lnot \lforall[{\OLId{latin-ordinary}{x}}][!A({\OLId{latin-ordinary}{x}})],
      \iftag{prvEx}{\lexists[{\OLId{latin-ordinary}{x}}][\lnot !A({\OLId{latin-ordinary}{x}})] \lif \lnot
        !A({\OLId{latin-ordinary}{c}})}{\lnot\lforall[{\OLId{latin-ordinary}{x}}][!A({\OLId{latin-ordinary}{x}})] \lif \lnot
        !A({\OLId{latin-ordinary}{c}})}, !A({\OLId{latin-ordinary}{c}})
      \]
      غير متسقة.)

      وأما الاتجاه المقابل فلا يفتقر إلى الإشباع. فإذا كان
      $\lforall[{\OLId{latin-ordinary}{x}}][!A({\OLId{latin-ordinary}{x}})] \in {\OLId{greek-variable}{Gamma}}$، لزم ${\OLId{greek-variable}{Gamma}} \Proves !A({\OLId{latin-ordinary}{t}})$
      بمقتضى
      \tagrefs{prfAX/{fol:axd:qpr:prop:provability-quantifiers},prfSC/{fol:seq:qpr:prop:provability-quantifiers},prfND/{fol:ntd:qpr:prop:provability-quantifiers},prfTab/{fol:tab:qpr:prop:provability-quantifiers}},
      البند~(\OLClassicalDigits{2}). ونحصل على $!A({\OLId{latin-ordinary}{t}}) \in {\OLId{greek-variable}{Gamma}}$ بحسب
      \olref[ccs]{prop:ccs}.}}{}
  \end{tagenumerate}
\end{proof}

\end{document}
